\section{Socles, Tops and Projective Covers}

\begin{definition}[Soc of a module]
    The unique largest semisimple submodule of an $R$-module $U$ is called the
    \emph{socle} of $U$ and is denoted by
    \[
        \operatorname{Soc}(U).
    \]
\end{definition}

\begin{definition}[Top of a module]
    Let $U$ be an $R$-module. The top, or head, of $U$ is defined by
    \[
        \operatorname{Top}(U):=U/\operatorname{Rad}(U).
    \]
\end{definition}

\begin{example}
    If $U$ is semisimple, then
    \[
        \operatorname{Rad}(U)=0,
        \qquad
        \operatorname{Soc}(U)=U,
        \qquad
        \operatorname{Top}(U)=U.
    \]
    In particular, if $S$ is simple, then
    \[
        \operatorname{Rad}(S)=0,
        \qquad
        \operatorname{Soc}(S)=S,
        \qquad
        \operatorname{Top}(S)=S.
    \]
\end{example}

\begin{example}
    Let
    \[
        A=k[x]/(x^n),
    \]
    and regard $A$ as a left $A$-module. Then
    \[
        J(A)=(x),
    \]
    so
    \[
        \operatorname{Rad}(A)=xA,
        \qquad
        \operatorname{Top}(A)=A/xA\cong k.
    \]
    On the other hand,
    \[
        \operatorname{Soc}(A)=\{a\in A\mid xa=0\}=kx^{n-1}.
    \]
    Thus
    \[
        \operatorname{Top}(A)\cong k,
        \qquad
        \operatorname{Soc}(A)\cong k,
    \]
    but they occur at opposite ends of the Loewy filtration.
\end{example}

\begin{remark}
    In general, $\operatorname{Top}(U)$ need not be semisimple for an arbitrary module
    $U$. However, if $U$ is finitely generated and Artinian, then
    $\operatorname{Top}(U)$ is semisimple.
\end{remark}

\begin{lemma}
    Let
    \[
        U=S_1^{\oplus a_1}\oplus \cdots \oplus S_r^{\oplus a_r}
    \]
    be a semisimple \(R\)-module, where the \(S_i\) are pairwise nonisomorphic
    simple modules. Then the summand \(S_i^{\oplus a_i}\) is the \(S_i\)-isotypic
    component of \(U\), namely the sum of all submodules of \(U\) isomorphic to
    \(S_i\). In particular, it is uniquely determined by \(U\) and the isomorphism
    class of \(S_i\).
\end{lemma}

\begin{proof}
    This is precisely Corollary~5.2.5(2) applied to \(S=S_i\).
\end{proof}

\begin{lemma}
    Let $U$ be an $R$-module.
    \begin{enumerate}
        \item Suppose $M_1,\dots,M_n$ are maximal submodules of $U$. Then there exists a
              subset $I\subseteq \{1,\dots,n\}$ such that
              \[
                  U/(M_1\cap \cdots \cap M_n)\cong \bigoplus_{i\in I} U/M_i.
              \]
              In particular, $U/(M_1\cap \cdots \cap M_n)$ is semisimple.
        \item If $U$ is finitely generated and Artinian, then $U/\operatorname{Rad}(U)$ is
              semisimple, and $\operatorname{Rad}(U)$ is the unique smallest submodule
              of $U$ with semisimple quotient.
    \end{enumerate}
\end{lemma}

\begin{proof}
    \textbf{(1).}
    Choose a subset $I\subseteq \{1,\dots,n\}$ maximal with the property that the
    diagonal map
    \[
        \delta_I:U\longrightarrow \bigoplus_{i\in I} U/M_i,\qquad
        u\longmapsto (u+M_i)_{i\in I}
    \]
    is surjective. Such a subset exists because $\{1,\dots,n\}$ is finite; for instance
    $I=\varnothing$ is allowed.

    We have
    \[
        \ker(\delta_I)=\bigcap_{i\in I}M_i.
    \]
    Put
    \[
        K:=\bigcap_{i\in I}M_i.
    \]
    We claim that
    \[
        K\subseteq M_j
        \qquad\text{for every }j\notin I.
    \]
    Suppose not. Then for some $j\notin I$ we have
    \[
        K\nsubseteq M_j.
    \]
    Since $M_j$ is maximal, this implies
    \[
        K+M_j=U.
    \]
    Hence the natural map
    \[
        K\longrightarrow U/M_j
    \]
    is surjective.

    We now show that
    \[
        U\longrightarrow
        \left(\bigoplus_{i\in I}U/M_i\right)\oplus U/M_j
    \]
    is surjective. Indeed, given
    \[
        \bigl((u_i+M_i)_{i\in I},\, u_j+M_j\bigr),
    \]
    first use the surjectivity of $\delta_I$ to choose $u\in U$ such that
    \[
        u+M_i=u_i+M_i
        \qquad\text{for all }i\in I.
    \]
    Since $K\to U/M_j$ is surjective, choose $k\in K$ such that
    \[
        k+M_j=(u_j-u)+M_j.
    \]
    Then $u+k$ maps to the prescribed element in every component:
    for $i\in I$ we have $k\in K\subseteq M_i$, so
    \[
        u+k+M_i=u+M_i=u_i+M_i,
    \]
    and for the $j$-component,
    \[
        u+k+M_j=u_j+M_j.
    \]
    This contradicts the maximality of $I$.

    Therefore
    \[
        \bigcap_{i\in I}M_i\subseteq M_j
        \qquad\text{for every }j\notin I.
    \]
    Hence
    \[
        \bigcap_{i\in I}M_i
        =
        M_1\cap\cdots\cap M_n.
    \]
    By the first isomorphism theorem,
    \[
        U/(M_1\cap\cdots\cap M_n)
        =
        U/\ker(\delta_I)
        \cong
        \bigoplus_{i\in I}U/M_i.
    \]
    Since each $M_i$ is maximal, each quotient $U/M_i$ is simple. Therefore the right
    hand side is semisimple.

    \textbf{(2).}
    If $U=0$, the statement is trivial. Assume $U\neq 0$.

    Since $U$ is finitely generated, every proper submodule of $U$ is contained in a
    maximal submodule. Indeed, if $W\subsetneq U$, then $U/W$ is finitely generated,
    so by the maximal-submodule lemma there exists a maximal submodule $M$ of $U$
    such that
    \[
        W\subseteq M\subsetneq U.
    \]

    We now show that the intersection of all maximal submodules is in fact a finite
    intersection. Consider the set
    \[
        \mathcal{F}
        :=
        \left\{
            M_1\cap\cdots\cap M_t
            \ \middle|\
            t\geq 1,\ M_1,\dots,M_t\text{ maximal submodules of }U
        \right\}.
    \]
    This is a nonempty set of submodules of $U$. Since $U$ is Artinian, every nonempty
    set of submodules has a minimal element. Choose a minimal element
    \[
        K=M_1\cap\cdots\cap M_t
        \]
    in $\mathcal{F}$.

    Let $M$ be any maximal submodule of $U$. Then
    \[
        K\cap M=M_1\cap\cdots\cap M_t\cap M
    \]
    also belongs to $\mathcal{F}$, and
    \[
        K\cap M\subseteq K.
    \]
    By the minimality of $K$, we must have
    \[
        K\cap M=K.
    \]
    Hence
    \[
        K\subseteq M.
    \]
    Since this holds for every maximal submodule $M$ of $U$, we get
    \[
        K\subseteq \operatorname{Rad}(U).
    \]
    The opposite inclusion is immediate because $K$ is an intersection of maximal
    submodules. Therefore
    \[
        \operatorname{Rad}(U)=K=M_1\cap\cdots\cap M_t.
    \]

    By part (1), it follows that
    \[
        U/\operatorname{Rad}(U)
    \]
    is semisimple.

    It remains to prove the minimality statement. Let $V\leq U$ be a submodule such
    that $U/V$ is semisimple. We show that
    \[
        \operatorname{Rad}(U)\subseteq V.
    \]
    Since $U/V$ is semisimple, its radical is zero:
    \[
        \operatorname{Rad}(U/V)=0.
    \]
    Maximal submodules of $U/V$ correspond exactly to maximal submodules of $U$
    containing $V$. Therefore
    \[
        \operatorname{Rad}(U/V)
        =
        \bigcap_{\substack{M\text{ maximal in }U\\ V\subseteq M}}
        M/V.
    \]
    Thus
    \[
        0
        =
        \operatorname{Rad}(U/V)
        =
        \left(
        \bigcap_{\substack{M\text{ maximal in }U\\ V\subseteq M}}
        M
        \right)/V.
    \]
    Hence
    \[
        \bigcap_{\substack{M\text{ maximal in }U\\ V\subseteq M}}M
        =
        V.
    \]
    Since $\operatorname{Rad}(U)$ is the intersection of all maximal submodules of $U$,
    it is contained in the intersection over the smaller family of maximal submodules
    containing $V$. Therefore
    \[
        \operatorname{Rad}(U)\subseteq V.
    \]

    Hence $\operatorname{Rad}(U)$ is contained in every submodule $V$ such that
    $U/V$ is semisimple. Since we already proved that $U/\operatorname{Rad}(U)$ is
    semisimple, $\operatorname{Rad}(U)$ is the unique smallest submodule of $U$ with
    semisimple quotient.
\end{proof}

\begin{corollary}[Universal property of the top]
    Let $U$ be a finitely generated Artinian $R$-module. Then
    $\operatorname{Top}(U)$ is the largest semisimple quotient of $U$ in the following sense:
    if $S$ is semisimple and
    \[
        f:U\to S
    \]
    is an $R$-module homomorphism, then $f$ factors uniquely through the canonical
    quotient map
    \[
        \pi:U\to \operatorname{Top}(U).
    \]
    Equivalently, there exists a unique homomorphism
    \[
        \overline{f}:\operatorname{Top}(U)\to S
    \]
    such that
    \[
        f=\overline{f}\circ \pi.
    \]
\end{corollary}

\begin{example}
    Suppose $U$ has finite length. Since $\operatorname{Top}(U)$ is semisimple, we may write
    \[
        \operatorname{Top}(U)
        \cong
        S_1^{\oplus a_1}\oplus \cdots \oplus S_r^{\oplus a_r},
    \]
    where the $S_i$ are pairwise nonisomorphic simple modules. The integers $a_i$
    measure how many copies of $S_i$ occur in the top layer of $U$.
\end{example}

\begin{lemma}
    For all $R$-modules $U$ and $V$,
    \[
        \operatorname{Rad}(U\oplus V)
        =
        \operatorname{Rad}(U)\oplus \operatorname{Rad}(V),
    \]
    \[
        \operatorname{Soc}(U\oplus V)
        =
        \operatorname{Soc}(U)\oplus \operatorname{Soc}(V),
    \]
    \[
        \operatorname{Top}(U\oplus V)
        \cong
        \operatorname{Top}(U)\oplus \operatorname{Top}(V).
    \]
\end{lemma}

\begin{proof}
    The projection maps show that every maximal submodule of $U\oplus V$ is of the form
    \[
        M\oplus V
        \qquad\text{or}\qquad
        U\oplus N,
    \]
    where $M$ and $N$ are maximal submodules of $U$ and $V$, respectively. Intersecting
    all such maximal submodules yields
    \[
        \operatorname{Rad}(U\oplus V)=\operatorname{Rad}(U)\oplus \operatorname{Rad}(V).
    \]

    For the socle, every simple submodule of $U\oplus V$ projects nontrivially onto at
    least one of the two factors; by simplicity it must embed into one factor. Therefore
    every simple submodule of $U\oplus V$ lies in
    \[
        \operatorname{Soc}(U)\oplus \operatorname{Soc}(V),
    \]
    and the reverse inclusion is clear.
\end{proof}

\begin{proposition}
    Let
    \[
        f:P\to U
    \]
    be a projective cover, and assume that $P$ and $U$ are finitely generated Artinian.
    Then the induced morphism
    \[
        \overline{f}:\operatorname{Top}(P)\to \operatorname{Top}(U)
    \]
    is an isomorphism.
\end{proposition}

\begin{proposition}
    Suppose that
    \[
        f:P\to U
    \]
    is a projective cover of an $R$-module $U$, and that
    \[
        g:Q\to U
    \]
    is an epimorphism with $Q$ projective.
    \begin{enumerate}
        \item There exists a decomposition
              \[
                  Q=Q_1\oplus Q_2
              \]
              such that, with respect to this decomposition,
              \[
                  g=(g_1,0),
              \]
              and there is an isomorphism $\gamma:Q_1\to P$ satisfying
              \[
                  f\circ \gamma=g_1.
              \]
        \item If a projective cover of $U$ exists, then it is unique up to
              isomorphism.
    \end{enumerate}
\end{proposition}

\begin{proof}
    Since $Q$ is projective and $f:P\to U$ is surjective, there exists
    \[
        \alpha:Q\to P
    \]
    such that
    \[
        f\circ \alpha=g.
    \]
    Since $P$ is projective and $g:Q\to U$ is surjective, there exists
    \[
        \beta:P\to Q
    \]
    such that
    \[
        g\circ \beta=f.
    \]
    Consequently,
    \[
        f=f\circ (\alpha\beta).
    \]
    Because $f$ is an essential epimorphism, the endomorphism $\alpha\beta$ of $P$ must
    be surjective. Since $P$ is projective, $\alpha\beta$ splits. Let
    \[
        P=\ker(\alpha\beta)\oplus P'.
    \]
    If $\ker(\alpha\beta)\neq 0$, then $P'$ is a proper submodule of $P$ and
    \[
        f(P')=f(\alpha\beta(P'))=f(P)=U,
    \]
    contradicting the essentiality of $f$. Hence $\ker(\alpha\beta)=0$, so $\alpha\beta$ is an automorphism.

    Replacing $\beta$ by $\beta\circ (\alpha\beta)^{-1}$, we may assume that
    \[
        \alpha\beta=\operatorname{id}_P.
    \]
    Thus $\beta$ is a split monomorphism and
    \[
        Q=\beta(P)\oplus \ker(\alpha).
    \]
    Set
    \[
        Q_1:=\beta(P),
        \qquad
        Q_2:=\ker(\alpha),
    \]
    and let $\gamma:Q_1\to P$ be the inverse of the isomorphism $\beta:P\to Q_1$.
    Then
    \[
        g|_{Q_1}=f\circ \gamma.
    \]
    Moreover,
    \[
        g(Q_2)=f(\alpha(Q_2))=0,
    \]
    so $g=(g_1,0)$ with $g_1=f\circ \gamma$. This proves (1).

    For (2), if
    \[
        f:P\to U
        \qquad\text{and}\qquad
        f':P'\to U
    \]
    are projective covers, apply (1) twice: once with $Q=P'$ and once with $Q=P$. We
    deduce that $P$ is isomorphic to a direct summand of $P'$ and $P'$ is isomorphic to a
    direct summand of $P$. The decompositions obtained in (1) force the complementary
    summands to map trivially to $U$, hence to vanish by essentiality. Therefore
    \[
        P\cong P',
    \]
    so the projective cover is unique up to isomorphism.
\end{proof}

